\section*{Capítulo \ref{ch:g8:ohms-law} --- A resistência e a lei de Ohm}

\begin{solution}{exo:g8:ohms-law:1}
Com que \omterm{def:g2:pushes-and-pulls:force}{força} um componente se opõe à corrente; \omterm{def:g6:measurement-in-science:measuring}{unidade}, o \omterm{def:g8:ohms-law:resistance}{ohm};
símbolo, $\Omega$. $U = R I$; \ $I = U/R$; \ $R = U/I$.
\end{solution}

\begin{solution}{exo:g8:ohms-law:2}
$I = 6 \div 30 = \qty{0.2}{A}$; \ $U = 30 \times 0.5 = \qty{15}{V}$.
\end{solution}

\begin{solution}{exo:g8:ohms-law:3}
$R = 4.5 \div 0.09 = 50\,\Omega$.
\end{solution}

\begin{solution}{exo:g8:ohms-law:4}
Uma reta que passa pela origem --- a proporcionalidade. Uma reta mais
inclinada quer dizer mais \omterm{def:g8:voltage-voltmeter:voltage}{volts} necessários por \omterm{def:g8:intensity-ammeter:intensity}{ampère}: \omterm{def:g8:ohms-law:resistance}{resistência}
maior.
\end{solution}

\begin{solution}{exo:g8:ohms-law:5}
$3.0 \div 0.20 = 15$; $6.0 \div 0.40 = 15$. A igualdade da razão de
uma coluna para a outra \emph{é} a lei: um único $R$ constante
servindo a cada leitura é o que faz de ``a \omterm{def:g8:ohms-law:resistance}{resistência}'' uma
propriedade do componente.
\end{solution}

\begin{solution}{exo:g8:ohms-law:6}
A \omterm{def:g8:ohms-law:resistance}{resistência} menor bebe mais: $230 \div 25 = \qty{9.2}{A}$ contra
$230 \div 35 \approx \qty{6.6}{A}$.
\end{solution}

\begin{solution}{exo:g8:ohms-law:7}
Lâmpadas \omterm{def:g4:series-circuits:series}{em série} repartem um mesmo $I$; a \omterm{def:g8:voltage-voltmeter:voltage}{tensão} de cada lâmpada é
$U = R I$ com o mesmo $I$ --- então o maior $R$ multiplica até o maior
$U$: quem se opõe mais fica com a parcela maior.
\end{solution}

\begin{solution}{exo:g8:ohms-law:8}
\omterm{def:g7:short-circuits-safety:device}{Curto-circuito}: $R$ perto de zero, e então $I = U/R$ explode --- o
estouro da boiada sempre foi uma divisão por quase nada. \omterm{ex:g3:first-electric-circuit:switch}{Interruptor}
aberto: $R$ além de qualquer medida, e então $I = U/R$ desaba para
zero --- o bloqueio total.
\end{solution}

\begin{solution}{exo:g8:ohms-law:9}
O \omterm{def:g8:ohms-law:resistor}{resistor} precisa ficar com $9 - 2 = \qty{7}{V}$ a \qty{0.02}{A}:
$R = 7 \div 0.02 = 350\,\Omega$.
\end{solution}

\begin{solution}{exo:g8:ohms-law:10}
Primeiro ponto: $R = 1 \div 0.25 = 4\,\Omega$; segundo:
$4 \div 0.5 = 8\,\Omega$ --- valor único nenhum serve. A história: a
corrente \omterm{def:g7:moon-phases:phases}{crescente} aquece o filamento, e metal quente resiste mais; o
retrato entorta porque o componente muda enquanto trabalha.
\end{solution}

\begin{solution}{exo:g8:ohms-law:11}
Uma volta, um $I$: as duas quedas somam o empurrão da bateria,
$6 = 10 I + 20 I = 30 I$, então $I = \qty{0.2}{A}$. Daí
$U_{10} = 10 \times 0.2 = \qty{2}{V}$ e
$U_{20} = 20 \times 0.2 = \qty{4}{V}$ --- somando \qty{6}{V}, como a
lei da volta exige.
\end{solution}

\begin{solution}{exo:g8:ohms-law:12}
Cada \omterm{def:g5:series-parallel-circuits:parallel}{ramo} se estende sobre \qty{6}{V}:
$I_{10} = 6 \div 10 = \qty{0.6}{A}$,
$I_{20} = 6 \div 20 = \qty{0.3}{A}$; total \qty{0.9}{A}. A
substituta: $R = 6 \div 0.9 \approx 6.7\,\Omega$ --- \emph{menos} do
que qualquer um dos membros. Na linguagem das estradas: duas estradas
entre as mesmas cidades carregam mais tráfego do que qualquer uma
sozinha; toda estrada paralela a mais facilita a passagem, então a
oposição da dupla cai abaixo da do membro mais fraco.
\end{solution}

\begin{solution}{pb:g8:ohms-law:1}
\textbf{1.} Fonte, \omterm{def:g8:intensity-ammeter:ammeter}{amperímetro} e componente numa volta \omterm{def:g4:series-circuits:series}{em série} ---
o pedágio dentro da estrada; o \omterm{def:g8:voltage-voltmeter:voltmeter}{voltímetro} atravessado no componente
--- o topógrafo de lado.
\textbf{2.} $R = 4.5 \div 0.15 = 30\,\Omega$.
\textbf{3.} Cinco por cento de $33$ dão cerca de $1.65$: aceitável de
$31.35$ a $34.65\,\Omega$. Os $30\,\Omega$ da amostra um caem fora:
rejeitada.
\textbf{4.} $I = 4.5 \div 150 = \qty{0.03}{A} = \qty{30.0}{mA}$.
\textbf{5.} Toda coluna devolve $R = U/I = 30\,\Omega$
($1.5 \div 0.05$, $3.0 \div 0.10$, \dots): razão constante, retrato
em reta --- um valor serve para tudo: $30\,\Omega$.
\textbf{6.} As razões dão $5$, $6.7$, $8.3$, $10\,\Omega$: subindo a
cada passo --- \omterm{def:g8:ohms-law:resistance}{resistência} crescendo conforme a corrente cresce;
nenhum $R$ único existe.
\textbf{7.} Conforme a corrente cresce, o filamento esquenta em
direção ao brilho, e metal quente resiste mais --- o retrato entorta
para cima. O teste da reta da bancada admite portanto exatamente os
componentes cujo $R$ fica parado: ele é, por construção, um detector
de \omterm{def:g8:ohms-law:resistor}{resistores}.
\textbf{8.} A cláusula da \omterm{def:g3:temperature-thermometers:temperature}{temperatura} estável: a lei de Ohm promete
proporcionalidade \emph{a \omterm{def:g3:temperature-thermometers:temperature}{temperatura} estável} --- esquente um
componente no meio do retrato e até um \omterm{def:g8:ohms-law:resistor}{resistor} honesto deriva.
\textbf{9.} $R = 12 \div 0.03 = 400\,\Omega$.
\textbf{10.} $U = R I = 25 \times 9.2 = \qty{230}{V}$ --- exatamente
a tomada: a \omterm{def:g8:ohms-law:resistance}{resistência} puxa precisamente o que a lei de Ohm lhe
destina, e a margem do \omterm{def:g7:short-circuits-safety:fuse}{fusível} funciona como foi projetada. Nada é
``só''; tudo é aritmética.
\textbf{11.} O fio nasceu para ser um fio de ligação --- a estrada
plana, que não gasta empurrão. Mal usado, atravessado numa fonte, os
\omterm{def:g8:ohms-law:resistance}{ohms} quase nulos dele fazem o \omterm{def:g7:short-circuits-safety:device}{curto-circuito} perfeito: o vilão sempre
foi apenas um fio muito bom no lugar errado.
\textbf{12.} Por exemplo: ``A lei: $U = R I$, empurrão igual a
oposição vezes marcha. O retrato dela: uma reta pela origem, mais
inclinada para quem se opõe mais. E só um componente a \omterm{def:g3:temperature-thermometers:temperature}{temperatura}
estável pode assiná-la --- o \omterm{def:g5:heat-and-insulation:heat}{calor} reescreve a \omterm{def:g8:ohms-law:resistance}{resistência}, e
retratos tirados com febre não provam nada.''
\end{solution}
